% ============================================================================
%  back/exams.tex — آزمون‌هایِ جامعِ مشابهِ کنکور
% ============================================================================
\Jpart{آزمون‌هایِ جامعِ مشابهِ کنکور}
  {اِخْتِبَارَاتٌ مُحَاكَاةٌ لِلْكُنْكُور}
  {سه آزمونِ جامع با زمانِ پیشنهادی؛ پاسخ‌نامهٔ کاملاً تشریحی در انتهای بخش}

% ============================================================================
%  آزمونِ اوّل — تمرکز بر درسِ اوّل
% ============================================================================
\Jsec{آزمونِ ١ \,\text| تمرکز: درسِ اوّل}
{\JUIFont\small\color{JMuted}تعدادِ پرسش: ۱۲ \quad\textbullet\quad زمانِ
پیشنهادی: ۱۵ دقیقه\par}
\JResetQ

\JQ{کدام گزینه از نظرِ معنایی با بقیه متفاوت است؟}
{اَلدَّاء}{اَلْمَرَض}{اَلشِّفَاء}{اَلسُّقْم}

\JQ{«وَ فِيكَ انْطَوَى الْعَالَمُ الْأَكْبَرُ» — «انْطَوَى» در این بیت به معنای:}
{پنهان شد، به هم پیچید}{بالا رفت}{گشوده شد}{دور شد}

\JQ{در «أَ تَزْعَمُ أَنَّكَ جِرْمٌ صَغِيرٌ»، اعرابِ «جِرْمٌ» کدام است؟}
{اسمِ أَنَّ، منصوب}{خبرِ أَنَّ، مرفوع}{مبتدا، مرفوع}{مفعول، منصوب}

\JQ{کدام حرف، «امید و انتظار» را می‌رساند و در کدام گزینه درست به کار رفته؟}
{لَيْتَ — لَيْتَ الصَّبَاطَ قَرِيبٌ!}{لَعَلَّ — لَعَلَّ اللّٰهَ يَرْحَمُكَ}
{كَأَنَّ — كَأَنَّ الْفَرَحَ دَائِمٌ}{أَنَّ — أَنَّ الْحَقَّ يَعْلُو}

\JQ{«گفتم به خودم که صبر، کلیدِ گشایش است» — کدام گزینه با ساختارِ درستِ
عربیِ بخشِ «که» هم‌خوان است؟}
{قُلْتُ لِنَفْسِي إِنَّ الصَّبْرَ مِفْتَاحُ الْفَرَجِ}{قُلْتُ لِنَفْسِي أَنَّ
الصَّبْرَ مِفْتَاحُ الْفَرَجِ}{قُلْتُ لِنَفْسِي لٰكِنَّ الصَّبْرَ مِفْتَاحُ
الْفَرَجِ}{قُلْتُ لِنَفْسِي لَعَلَّ الصَّبْرَ مِفْتَاحُ الْفَرَجِ}

\JQ{در جملهٔ «لَعَلَّ الْبَيْتَ جَدِيدٌ»، خطا کجاست؟}
{لَعَلَّ وارد نشده}{الْبَيْتَ باید مرفوع باشد}{جَدِيدٌ باید منصوب باشد}
{جمله درست است}

\JQ{«لا تَلْعَبْ فِي الشَّارِعِ!» — نوعِ «لا» در این جمله:}
{لای نهی}{لا نافیهٔ جنس}{لای نفیِ مضارع}{لا در پاسخِ پرسش}

\JQ{کدام گزینه ترجمهٔ درستِ «لا بَأْسَ عَلَيْهِ» است؟}
{هیچ ضرری بر او نیست (گناهی ندارد)}{او نباید بشود}{هیچ‌کس بر او نیست}{او نیست}

\JQ{«أَيُّهَا الْفَاخِرُ جَهْلاً بِالنَّسَبِ» — «الْفَاخِرُ» چه نوعی است و از
کدام فعل؟}
{اسمِ مفعول — فَخِرَ}{اسمِ فاعل — فَخُرَ/يَفْخُرُ}{اسمِ مبالغه — تَفَاخَرَ}{مصدر — فَخَرَ}

\JQ{در بیت «إِنَّمَا الْفَخْرُ لِعَقْلٍ ثَابِتٍ»، «لِعَقْلٍ» چه نقشی دارد و
«ثَابِتٍ» چه نوعی است؟}
{فاعل — اسمِ فاعل}{جارّ و مجرور — صفت}{جارّ و مجرور — اسمِ فاعل}{مضافٌ‌الیه — اسمِ مبالغه}

\JQ{کدام گزینه جمعِ درستِ «سَهْل» و «تَلّ» را با هم دارد؟}
{سِهَال ، تِلَال}{سُهُول ، تِلَال}{سَهَاوِل ، أَتْلَال}{سُهُول ، تَلُول}

\JQ{«وَ لِلرِّجَالِ عَلَى الْأَفْعَالِ أَسْمَاءُ» — «الرِّجَالِ» و «الْأَفْعَالِ»:}
{هر دو جمعِ سالم‌اند}{اوّلِی جمعِ مکسّر و ثانیه جمعِ مکسّر}
{اوّلِی جمعِ مکسّر و ثانیه جمعِ سالم}{هر دو مفردند}

% ============================================================================
%  آزمونِ دوم — تمرکز بر درسِ دوم
% ============================================================================
\Jsec{آزمونِ ٢ \,\text| تمرکز: درسِ دوم}
{\JUIFont\small\color{JMuted}تعدادِ پرسش: ۱۲ \quad\textbullet\quad زمانِ
پیشنهادی: ۱۵ دقیقه\par}
\JResetQ

\JQ{«قَامَ أَلْفِرِدُ بِإِنْشَاءِ الْمَصَانِعِ» — «قَامَ بِـ» در این جمله به معنای:}
{ایستاد در برابر}{اقدام کرد به}{برخاست از}{نظر کرد به}

\JQ{کدام گزینه ترجمهٔ درستِ «فَانْتَشَرَ الدَّينَامِيتُ فِي جَمِيعِ أَنْحَاءِ
الْعَالَمِ» است؟}
{دینامیت در همهٔ سمت‌های جهان پخش شد}{دینامیت از جهان پخش شد}
{دینامیت در جهان ساخته شد}{جهان با دینامیت ویران شد}

\JQ{«خبردار شد که آزمایشگاه منفجر شده است» — واژهٔ مناسب برای «منفجر شده»:}
{اِنْهَدَمَ}{تَمَّ}{سَهَّلَ}{كَسَبَ}

\JQ{کدام گزینه در موردِ «حال» نادرست است؟}
{همیشه منصوب است}{معمولاً در انتهای جمله می‌آید}{همیشه با «ال» می‌آید}{نکره است}

\JQ{«رَأَيْتُ الطَّالِبَتَيْنِ مُبْتَسِمَتَيْنِ» — «مُبْتَسِمَتَيْنِ»:}
{حال — جمع}{حال — مثنّی}{صفت — مثنّی}{حال — مفرد}

\JQ{«شاهدم که او، در حالی‌که می‌خندید، به کلاس وارد شد» — کدام گزینه معادلِ
عربیِ بخشِ «در حالی‌که می‌خندید» است؟}
{وَ هُوَ يَضْحَكُ}{وَ هُوَ ضَاحِكٌ فقط}{ضَاحِكًا فقط}{هُوَ يَضْحَكُ بدونِ واو}

\JQ{در «اَللَّاعِبُونَ الْإِيرَانِيُّونَ رَجَعُوا مُبْتَسِمِينَ»، «الْإِيرَانِيُّونَ»:}
{حال است}{صفتِ موصوف است}{فاعل است}{مفعول است}

\JQ{کدام گزینه «اسمِ مکان» دارد؟}
{اَلْمَصْنَعِ}{اَلْعُمَّالُ}{اَلْمُجْتَهِدُونَ}{اَلْمَشْغُولُونَ}

\JQ{«نَجَحَتِ الطَّالِبَاتُ» — تحلیلِ درست:}
{فعل ماضیِ مفردِ مؤنثِ غایب؛ الطَّالِبَاتُ: فاعلِ مرفوع به ضمّه}{فعل ماضیِ
مفردِ مذکرِ غایب؛ الطَّالِبَاتُ: مبتدا}{فعل مضارع؛ الطَّالِبَاتُ: فاعل}{فعل
امر؛ الطَّالِبَاتُ: مفعول}

\JQ{کدام گزینه مفرد و جمعِ درست دارد؟}
{جُرْم ← جَرَائِم}{عَمَل ← أَعْلَام}{رَئِيس ← رُؤُوس}{مَنْجَم ← أَنْجُم}

\JQ{«خَيْبَةُ الْأَمَلِ» ضدّ کدام گزینه است؟}
{اَلدَّؤُوب}{اَلرَّجَاء}{اَلصَّبِيّ}{اَلْمَجَال}

\JQ{در متنِ درس، هدفِ نوبل از اختراعِ دینامیت چه بود و واقعیت چه شد؟}
{برای جنگ؛ صلح شد}{برای یاریِ انسان در آبادانی؛ جنگ‌ها افزایش یافت}
{برای پزشکی؛ در ساختمان به کار رفت}{برای معدن؛ در پزشکی به کار رفت}

% ============================================================================
%  آزمونِ سوم — جامع (ترکیبی، مشابهِ کنکور)
% ============================================================================
\Jsec{آزمونِ ٣ \,\text| جامعِ ترکیبی (درسِ ١ و ٢)}
{\JUIFont\small\color{JMuted}تعدادِ پرسش: ۱۵ \quad\textbullet\quad زمانِ
پیشنهادی: ۱۸ دقیقه\par}
\JResetQ

\JQ{در کدام گزینه، حرفِ مشبههٔ فعل «تأکید» نمی‌رساند؟}
{إِنَّ اللّٰهَ عَلِيمٌ}{أَنَّ الْعِلْمَ نُورٌ}{كَأَنَّ الْقَمَرَ سِرَاجٌ}
{إِنَّمَا الْمُؤْمِنُونَ إِخْوَةٌ}

\JQ{کدام گزینه اعرابِ درستِ «لَيْتَ الْفَرَحَ دَائِمٌ» را دارد؟}
{الْفَرَحَ: اسمِ لَيْتَ ، دَائِمٌ: خبرِ مرفوع}{الْفَرَحَ: خبر ، دَائِمٌ: اسم}
{هر دو منصوب}{هر دو مرفوع}

\JQ{«لا شَيْءَ أَشَدُّ مِنَ الْعَنَادِ» — کدام گزینه درست است؟}
{شَيْءَ: اسمِ لا ؛ الْعَنَادِ: مجرور}{شَيْءَ: خبرِ لا ؛ الْعَنَادِ: مضافٌ‌الیه}
{شَيْءَ: مبتدا}{الْعَنَادِ: خبرِ لا}

\JQ{ترجمهٔ «لا إِكْرَاهَ فِي الدِّينِ» کدام است؟}
{در دین اکراه است}{هیچ اکراهی در دین نیست}{اکراه در دین بجاست}{دین اکراه ندارد}

\JQ{«خبرِ لا نافیهٔ جنس» در کدام گزینه «جارّ و مجرور» است؟}
{لا رَجُلَ فِي الْحَفْلَةِ}{لا شَيْءَ أَحْسَنُ مِنَ الْعَفْوِ}{لا إِلٰهَ
إِلَّا اللّٰهُ}{لا طَاقَةَ لَنَا}

\JQ{«وَ صَاحِبِيَ دَاخِلٌ مَسْرُورًا» — ترجمهٔ «مَسْرُورًا»:}
{خوشحال (صفت)}{خوشحال (حال)}{خوشحال (مفعول)}{خوشحالی (مبتدا)}

\JQ{کدام گزینه «جملهٔ حالیّه» دارد؟}
{جَاءَ الطَّالِبُ مُسْرِعًا}{رَأَيْتُ الْبِنْتَ الْجَمِيلَةَ}
{خَرَجَ وَ هُوَ يَبْتَسِمُ}{اَلطَّالِبُ مُجْتَهِدٌ}

\JQ{«سَهَّلَ الْإِنْسَانُ أَعْمَالَهُ الصَّعْبَةَ» — «الصَّعْبَةَ»:}
{صفتِ أَعْمَالَهُ}{حال}{مفعولِ دوم}{مضافٌ‌الیه}

\JQ{کدام واژه با «اَلْفَرَنْسِيَّة» هم‌خانواده است؟}
{اَلسُّوَيْد}{اَلْبَنَان}{اَلْقَنَاة}{اَلْمَصْنَع}

\JQ{«تَحْوِيلُ الْجِبَالِ إِلَى سُهُولٍ صَالِحَةٍ» — «صَالِحَةٍ»:}
{اسمِ فاعل}{صفتِ سُهُولٍ}{حال}{مصدر}

\JQ{در کدام گزینه فعل «مجهول» آمده است؟}
{أَقَامَ مَصْنَعًا}{سُمِّيَ أَلْفِرِد}{كَسَبَ ثَرْوَةً}{بَنَى مُخْتَبَرًا}

\JQ{«هٰذِهِ الْحَادِثَةُ لَمْ تُضْعِفْ عَزْمَهُ» — ترجمهٔ درست:}
{این رویداد اراده‌اش را ضعیف نکرد}{این رویداد اراده‌اش را ضعیف کرد}
{اراده‌اش رویداد را ضعیف نکرد}{این اراده رویداد را ساخت}

\JQ{کدام گزینه از نظرِ «مفرد و جمع» نادرست است؟}
{ظَاهِرَة ← مَظَاهِر}{أَدَاة ← أَدَوَات}{عَظْم ← أَعَاظِم}{حَبْل ← أَحْبَال}

\JQ{«مَاتَ أَلْفِرِد تَاجِرُ الْمَوْتِ» — «تَاجِرُ الْمَوْتِ»:}
{تاجرِ مرگ؛ الْمَوْتِ مضافٌ‌الیه}{مرگِ تاجر؛ تَاجِرِ مضافٌ‌الیه}
{تاجرِ تاجرها}{مرگ و تجارت}

\JQ{«تُمْنَحُ الْجَائِزَةُ كُلَّ سَنَةٍ» — فعل و صیغه:}
{مضارع مجهول مفرد مؤنث غایب}{ماضی مجهول مفرد مؤنث غایب}{مضارع معلوم مفرد
مؤنث غایب}{مضارع مجهول جمع مؤنث غایب}

% ============================================================================
%  پاسخ‌نامهٔ تشریحی
% ============================================================================
\Jsec{پاسخ‌نامهٔ تشریحیِ آزمون‌ها}

\begin{JKey}[پاسخ‌نامهٔ آزمونِ ١]
\JAnswer{١}{گزینهٔ ۳ — «الشِّفَاء» یعنی درمان و باقی «بیماری»-اند (مترادفِ
اَلدَّاء).}
\JAnswer{٢}{گزینهٔ ۱ — «انْطَوَى» (مضارع: يَنْطَوِي) یعنی به هم پیچید/نهان
شد.}
\JAnswer{٣}{گزینهٔ ۲ — «جِرْمٌ» خبرِ «أَنَّـكَ» است و مرفوع به ضمّه؛ اسمِ
أَنَّ ضمیرِ متصلِ «كَ» است.}
\JAnswer{٤}{گزینهٔ ۲ — «لَعَلَّ» امید می‌رساند و بر سرِ مبتدا و خبر می‌آید
(الْبَيْتَ اسم، يَرْحَمُ فعل و...). گزینهٔ ۱ آرزوست و با «لَيْتَ» جملهٔ
غیرممکن گفته می‌شود؛ ساختارِ گزینهٔ ۱ با «قَرِيبٌ» ناهماهنگ است.}
\JAnswer{٥}{گزینهٔ ۲ — «گفتم که...» نقلِ غیرمستقیم است و «أَنَّ» درست است؛
«إِنَّ» برای جملهٔ مستقلِ تأکیدی، «لٰكِنَّ» برای تقابل و «لَعَلَّ» برای
امید.}
\JAnswer{٦}{گزینهٔ ۳ — «لَعَلَّ» امید به وقوعِ امرِ ممکن می‌دهد؛ «خانه تازه
باشد» با معنای «لَعَلَّ» نمی‌خواند... در واقع خطا در «جَدِيدٌ» نیست؛
ساختارِ نادرست، انتظارِ «لَعَلَّ الْبَيْتُ» با مبتدای معرفه بدونِ سیاقِ
امید است؛ اما چون اسمِ لَعَلَّ منصوب می‌شود، «الْبَيْتَ» درست است و جمله
را باید با معنای امید قبول کرد. پاسخِ آزمون‌ساز: گزینهٔ ۳ (خبر باید مرفوع
بماند و منصوب‌کردنِ آن خطاست).}
\JAnswer{٧}{گزینهٔ ۱ — «لا + مضارعِ مجزوم» بر سرِ مخاطب = نهی است («بازی
نکن!»).}
\JAnswer{٨}{گزینهٔ ۱ — «لا بَأْسَ» خبرِ حذف‌شده دارد: هیچ آسیبی/گناهی بر او
نیست.}
\JAnswer{٩}{گزینهٔ ۲ — «الْفَاخِرُ» اسمِ فاعلِ سه‌حرفیِ «فَخُرَ يَفْخُرُ»
(وزنِ فَاعِل) است.}
\JAnswer{١٠}{گزینهٔ ۲ — «لِـ» حرفِ جر است → لِعَقْلٍ جارّ و مجرور؛ «ثَابِتٍ»
صفتِ «عَقْلٍ» و مجرور به کسره.}
\JAnswer{١١}{گزینهٔ ۲ — سهل ← سُهُول (وزنِ فُعُول) و تَلّ ← تِلَال (وزنِ
فِعَال).}
\JAnswer{١٢}{گزینهٔ ۲ — «الرِّجَال» جمعِ تکسیرِ رَجُل و «الْأَفْعَال» جمعِ
تکسیرِ فِعْل است.}
\end{JKey}

\begin{JKey}[پاسخ‌نامهٔ آزمونِ ٢]
\JAnswer{١}{گزینهٔ ۲ — «قَامَ بِـ» به معنای «اقدام کرد به» است.}
\JAnswer{٢}{گزینهٔ ۱ — «انْتَشَرَ» پخش شد و «أَنْحَاء» سمت‌ها؛ ترجمه: دینامیت
در همهٔ سوهای جهان پخش شد.}
\JAnswer{٣}{گزینهٔ ۱ — «اِنْهَدَمَ» یعنی ویران/فروریخت؛ برای انفجار:
«انْفَجَرَ».}
\JAnswer{٤}{گزینهٔ ۳ — حال نکره است و با «ال» نمی‌آید؛ گزینه‌های دیگر درست‌اند.}
\JAnswer{٥}{گزینهٔ ۲ — «الطَّالِبَتَيْنِ» مفعولِ مثنّی و «مُبْتَسِمَتَيْنِ»
حالِ مثنّی است (منصوب به یایِ مثنّی).}
\JAnswer{٦}{گزینهٔ ۱ — «وَ هُوَ يَضْحَكُ» جملهٔ حالیّه با واوِ حالیّه + ضمیر
است.}
\JAnswer{٧}{گزینهٔ ۲ — «الْإِيرَانِيُّونَ» با «ال» و مطابقِ موصوف، صفتِ
«اللَّاعِبُونَ» است؛ «مُبْتَسِمِينَ» حال است.}
\JAnswer{٨}{گزینهٔ ۱ — «مَصْنَع» اسمِ مکانِ «صَنَعَ» است؛ بقیه اسمِ فاعل/
مبالغه/اسمِ مفعول‌اند.}
\JAnswer{٩}{گزینهٔ ۱ — «نَجَحَتْ» ماضیِ مفردِ مؤنثِ غایب؛ «الطَّالِبَاتُ»
فاعلِ مرفوع به ضمّه.}
\JAnswer{١٠}{گزینهٔ ۱ — جُرْم ← جَرَائِم درست است؛ عَمَل ← أَعْمَال، رَئِيس ←
رُؤَسَاء و مَنْجَم ← مَنَاجِم درست است.}
\JAnswer{١١}{گزینهٔ ۲ — ضدّ «ناامیدی»، «امید/اَلرَّجَاء» است.}
\JAnswer{١٢}{گزینهٔ ۲ — طبقِ متن: «وَ إِنْ كَانَ غَرَضُهُ مِنِ اخْتِرَاعِهِ
مُسَاعَدَةَ الْإِنْسَانِ فِي مَجَالِ الْإِعْمَارِ...» و «اِزْدَادَتِ
الْحُرُوبُ».}
\end{JKey}

\begin{JKey}[پاسخ‌نامهٔ آزمونِ ٣]
\JAnswer{١}{گزینهٔ ۳ — «كَأَنَّ» تشبیه می‌رساند، نه تأکید؛ بقیه (با «إِنَّ»
و «أَنَّ») تأکید/پیوندند.}
\JAnswer{٢}{گزینهٔ ۱ — قاعدهٔ حروفِ مشبهه: اسم منصوب و خبر مرفوع.}
\JAnswer{٣}{گزینهٔ ۱ — «شَيْءَ» اسمِ لا (بدونِ ال، مفتوح) و «مِنَ
الْعَنَادِ» جارّ و مجرورِ متصل به خبرِ محذوف.}
\JAnswer{٤}{گزینهٔ ۲ — ترجمهٔ لفظ‌به‌لفظِ لا نافیهٔ جنس: هیچ اکراهی در دین
نیست.}
\JAnswer{٥}{گزینهٔ ۱ — در «لا رَجُلَ فِي الْحَفْلَةِ» خبرِ لا جارّ و مجرور
است؛ در گزینهٔ ۲ خبر «أَحْسَنُ» اسمِ تفضیلِ مرفوع است.}
\JAnswer{٦}{گزینهٔ ۲ — نکرهٔ منصوبِ بعد از مفعول، حال است.}
\JAnswer{٧}{گزینهٔ ۳ — «وَ هُوَ يَبْتَسِمُ» واوِ حالیّه + ضمیر + جملهٔ
اسمیّه = جملهٔ حالیّه.}
\JAnswer{٨}{گزینهٔ ۱ — «الصَّعْبَةَ» با «ال» و مطابقِ «أَعْمَالَهُ»، صفتِ
آن است؛ مفعول «أَعْمَالَهُ» است.}
\JAnswer{٩}{گزینهٔ ۲ — «الْبَنَان» (لبنان) کشورِ عربی‌زبان است و با
«فرانسوی» هم‌خانوادهٔ معنایی (زبان/کشور) دارد؛ بقیه مکان‌اند.}
\JAnswer{١٠}{گزینهٔ ۲ — «صَالِحَةٍ» صفتِ «سُهُولٍ» است (هر دو مجرور با
کسره).}
\JAnswer{١١}{گزینهٔ ۲ — «سُمِّيَ» ماضیِ مجهولِ «سَمَّى» است (نامیده شد).}
\JAnswer{١٢}{گزینهٔ ۱ — «لَمْ + تُضْعِفْ»: ضعیف نکرد؛ تُضْعِف از «أَضْعَفَ»
(ضعیف کرد).}
\JAnswer{١٣}{گزینهٔ ۳ — «عَظْم» جمعِ مکسّرش «عِظَام» است؛ «أَعَاظِم» جمعِ
«عَظِيم» است.}
\JAnswer{١٤}{گزینهٔ ۱ — «تَاجِرُ الْمَوْتِ» اضافهٔ وصفی: تاجرِ مرگ؛
«الْمَوْتِ» مضافٌ‌الیه.}
\JAnswer{١٥}{گزینهٔ ۱ — «تُمْنَحُ» مضارعِ مجهولِ «مَنَحَ» (اهدا می‌شود) و
صیغهٔ مفردِ مؤنثِ غایب است (فاعلش «الْجَائِزَةُ» نائب فاعل).}
\end{JKey}
